\subsection{参数盲谱支、二态闭包、半熵禁带与有限秩账本障碍}\label{subsec:conclusion-parameterblind-binaryclosure-halfentropy-ledger-obstruction}
在前文已经分别建立二维加权核的 \((2+8)\)-次分裂、bin-fold escort 温度族的末位 Bernoulli 塌缩、六倍窗 Fourier 零块、Lee--Yang 分歧链以及有限秩加法账本障碍之后，还可把这些看似异质的结论进一步压缩为一条更高阶的终端统一链。其核心现象是：矩阵侧的刚性偶宇称支对碰撞参数导数完全失明，信息侧的 escort 温度族在极限中闭合为单比特实验，算术侧的六倍窗不仅强制出现半熵级禁带，而且在 \(m\equiv 4\pmod 6\) 的无限子序列上同步压低碰撞、KL 缺口、Fourier 峰值与最大纤维超额，局部场侧的十次层与 Lee--Yang 链仅共享一个 \(3\)-主方向，而守恒侧则把实现维与结构维永久分叉。

\begin{theorem}[参数盲偶宇称支分解]\label{thm:conclusion-parameterblind-even-parity-splitting}
\leanverified{paper\_conclusion\_parameterblind\_even\_parity\_splitting}
沿用附录 \ref{app:real-input-40-zeta-uv} 的记号，定义
\[
\mathcal Q_8(z;u,v):=-P_8(z;u,v),
\qquad
\mathcal Q_8(0;u,v)=1,
\]
并令
\[
\zeta_{\mathrm{red}}(z;u,v):=\zeta_{\mathrm{core}}(z;u,v):=\mathcal Q_8(z;u,v)^{-1}.
\]
则在形式幂级数环中有精确分解
\[
\Delta(z;u,v)=(1-vz^2)\mathcal Q_8(z;u,v),
\qquad
\zeta(z;u,v)=\frac{1}{1-vz^2}\,\zeta_{\mathrm{red}}(z;u,v).
\]
并且：
\begin{enumerate}
  \item 对任意 \(r\ge 1\)，
  \[
  \partial_u^{\,r}\log\zeta(z;u,v)=\partial_u^{\,r}\log\zeta_{\mathrm{red}}(z;u,v).
  \]
  \item 对任意 \(n\ge 1\)，
  \[
  \Tr(M_{u,v}^n)=2v^{n/2}\mathbf 1_{\{2\mid n\}}+R_n(u,v),
  \]
  其中 \(R_n(u,v)\) 仅由至多 \(8\) 条其余非零谱支贡献。
  \item 因而 \(\pm\sqrt v\) 所对应的刚性偶宇称支对一切碰撞参数 \(u\) 的导数响应完全失明；全部 \(u\)-敏感谱信息都被压缩到八次因子 \(\mathcal Q_8\) 之内。
\end{enumerate}
\end{theorem}
\begin{proof}
由定理 \ref{thm:real-input-40-det-uv-2plus8}，
\[
\Delta(z;u,v)=(vz^2-1)P_8(z;u,v)=(1-vz^2)\mathcal Q_8(z;u,v),
\]
且 \(\mathcal Q_8(0;u,v)=1\)。故
\[
\zeta(z;u,v)=\frac{1}{1-vz^2}\,\zeta_{\mathrm{red}}(z;u,v).
\]
第一条立即来自
\[
\log\zeta(z;u,v)=-\log(1-vz^2)+\log\zeta_{\mathrm{red}}(z;u,v),
\]
因为第一项与 \(u\) 无关。

另一方面，推论 \ref{cor:real-input-40-zeta-uv-sqrtv-eigs} 已给出两条显式谱支
\[
\lambda_\pm=\pm\sqrt v.
\]
再由定理 \ref{thm:xi-time-part73-nonrigid-trace-newton-cayley-closure} 中
\[
R_n(u,v)=\Tr(M_{u,v}^n)-(\sqrt v)^n-(-\sqrt v)^n
\]
的定义，即得
\[
\Tr(M_{u,v}^n)=2v^{n/2}\mathbf 1_{\{2\mid n\}}+R_n(u,v),
\]
且 \(R_n\) 仅由 \(\mathcal Q_8\) 所承载的至多 \(8\) 条非刚性谱支贡献。第三条只是前两条的合并解释：\((1-vz^2)^{-1}\) 仅产生一条 rigid 的偶长度 primitive ladder，并在一切 \(u\)-导数下被完全消去。证毕。
\end{proof}

\begin{corollary}[碰撞响应的谱支撑定理]\label{cor:conclusion-collision-response-spectral-support}
\leanverified{paper\_conclusion\_collision\_response\_spectral\_support}
对任意 \(r\ge 1\) 定义
\[
\Xi_r(z;u,v):=\partial_u^{\,r}\log\zeta(z;u,v).
\]
则 \(\Xi_r\) 的全部极点都落在 \(\mathcal Q_8(z;u,v)=0\) 的零点上；而
\[
z=\pm v^{-1/2}
\]
对 \(\Xi_r\) 均为可去奇点。等价地，普适偶宇称支虽决定一整条偶长度周期塔，却不贡献任何碰撞敏感临界点。
\end{corollary}
\begin{proof}
由定理 \ref{thm:conclusion-parameterblind-even-parity-splitting}，
\[
\Xi_r(z;u,v)=\partial_u^{\,r}\log\zeta_{\mathrm{red}}(z;u,v)
=-\partial_u^{\,r}\log\mathcal Q_8(z;u,v).
\]
因此其奇性只能来自 \(\mathcal Q_8(z;u,v)=0\) 的零点，而因子 \((1-vz^2)^{-1}\) 已被全部 \(u\)-导数消去。证毕。
\end{proof}

\begin{theorem}[escort 族的二态闭包]\label{thm:conclusion-escort-two-state-closure}
\leanverified{paper\_conclusion\_escort\_two\_state\_closure}
沿用 bin-fold escort 分布 \(\pi_{m,q}\) 的记号。对任意 \(p,q\ge 0\)，令
\[
\theta(q):=\frac{1}{1+\varphi^{q+1}},
\qquad
W_m^{(q)}\sim \pi_{m,q},
\]
并定义尺度化纤维权
\[
Z_m^{(q)}:=\left(\frac{\varphi}{2}\right)^m d_m^{\mathrm{bin}}\bigl(W_m^{(q)}\bigr).
\]
则有极限
\[
\lim_{m\to\infty}D_{\mathrm{KL}}(\pi_{m,q}\Vert \pi_{m,p})
=
\theta(q)\log\frac{\theta(q)}{\theta(p)}
+
\bigl(1-\theta(q)\bigr)\log\frac{1-\theta(q)}{1-\theta(p)}.
\]
同时
\[
Z_m^{(q)}\Rightarrow Z^{(q)},
\qquad
Z^{(q)}\in\{1,\varphi^{-1}\},
\]
且
\[
\mathbb P\bigl(Z^{(q)}=\varphi^{-1}\bigr)=\theta(q),
\qquad
\mathbb P\bigl(Z^{(q)}=1\bigr)=1-\theta(q).
\]
因此，bin-fold escort 温度族在信息几何极限中闭合为一个二态模型：全部参数 \(q\)-依赖最终仅由末位变量 \(w_m\in\{0,1\}\) 承载。
\end{theorem}
\begin{proof}
第一式是定理 \ref{thm:conclusion-binfold-escort-full-fdiv-collapse} 在
\[
\Phi(t)=t\log t
\]
处的特化；第二式则正是推论 \ref{cor:fold-bin-escort-two-scale-residual} 的重述。两者共同说明：无论考察温度间的相对熵几何，还是考察 escort 样本的尺度化纤维权，极限对象都压缩为同一 Bernoulli 两点机制。证毕。
\end{proof}

\begin{theorem}[六倍窗的半熵禁带]\label{thm:conclusion-half-entropy-forbidden-band}
\leanverified{paper\_conclusion\_half\_entropy\_forbidden\_band}
沿用六倍窗 Fourier 零块的记号，并记
\[
\mathcal Z_m:=\{k\in \ZZ/(F_{m+2}\ZZ):\ \widehat d_m(k)=0\}.
\]
若
\[
m+2\equiv 0\pmod 6,
\]
并令
\[
d:=\frac{m+2}{2},
\]
则有
\[
S_{F_d}(F_{m+2})\subseteq \mathcal Z_m,
\qquad
|\mathcal Z_m|\ge F_d=F_{(m+2)/2}.
\]
从而沿子序列 \(m\equiv 4\pmod 6\) 有
\[
\liminf_{\substack{m\to\infty\\ m\equiv 4\ (\mathrm{mod}\ 6)}}
\frac{1}{m}\log |\mathcal Z_m|
\ge \frac12\log\varphi.
\]
换言之，在该算术子序列上，被强制湮灭的频率集合以全系统熵率一半的速度增长。
\end{theorem}
\begin{proof}
推论 \ref{cor:fold-zero-half-index-multiple6} 在 \(m+2\equiv 0\pmod 6\) 时给出
\[
S_{F_d}(F_{m+2})\subseteq \mathcal Z_m,
\qquad
|\mathcal Z_m|\ge F_d.
\]
再由 Binet 公式
\[
F_n=\frac{\varphi^n}{\sqrt5}+O(\varphi^{-n}),
\]
可得
\[
\frac{1}{m}\log |\mathcal Z_m|
\ge
\frac{1}{m}\log F_{(m+2)/2}
=
\frac{m+2}{2m}\log\varphi+o(1).
\]
令 \(m\to\infty\) 且 \(m\equiv 4\pmod 6\)，便得到所示下界。证毕。
\end{proof}

\begin{theorem}[六周期窗口上的确定性三重下压律]\label{thm:conclusion-sixperiod-deterministic-triple-downshift}
\leanverified{paper\_conclusion\_sixperiod\_deterministic\_triple\_downshift}
记
\[
M:=F_{m+2},
\qquad
\overline d_m:=\frac{2^m}{F_{m+2}}.
\]
若
\[
m\equiv 4\pmod 6,
\]
则必有
\[
|\mathcal Z_m|\ge F_{(m+2)/2}.
\]
从而碰撞、相对熵、Fourier 峰值与最大纤维超额同时满足确定性改进界
\[
\mathsf{Col}_m\le 1-\frac{F_{(m+2)/2}}{F_{m+2}},
\]
\[
D_{\mathrm{KL}}(p_m\|u)\le \log\!\bigl(F_{m+2}-F_{(m+2)/2}\bigr),
\]
\[
S_m\le 2^m\sqrt{F_{m+2}-1-F_{(m+2)/2}},
\]
\[
M_m-\overline d_m\le
2^m\sqrt{\frac{F_{m+2}-1-F_{(m+2)/2}}{F_{m+2}}}.
\]
因此沿无限同余类 \(m\equiv 4\pmod 6\)，折叠频谱的碰撞能、信息缺口与峰值振幅会被同一条 Fibonacci 半指标零块同步压低，而最大纤维相对均值的凸起也随之进入相同的算术抑制轨道。
\end{theorem}
\begin{proof}
若 \(m\equiv 4\pmod 6\)，则 \(m+2\equiv 0\pmod 6\)。记
\[
d:=\frac{m+2}{2}.
\]
由推论 \ref{cor:fold-zero-half-index-multiple6}，
\[
|\mathcal Z_m|\ge F_d=F_{(m+2)/2}.
\]
再将这一半阶零块下界代入推论 \ref{cor:fold-zero-divisor-triple-reduction}，便得到四条确定性改进界。证毕。
\end{proof}

\begin{theorem}[共局部分歧的 \(3\)-主唯一性]\label{thm:conclusion-common-local-ramification-only-3}
\leanverified{paper\_conclusion\_common\_local\_ramification\_only\_3}
令 \(L_{10}\) 为十次层多项式 \(R_{10}(X)\) 的分裂域，\(L_{\mathrm{LY}}\) 为 Lee--Yang 三次数域 \(K=\QQ(y_{\mathrm{LY}})\) 的伽罗瓦闭包。记
\[
\mathcal R_{10}:=\mathrm{Ram}(L_{10}/\QQ),
\qquad
\mathcal R_{\mathrm{LY}}:=\mathrm{Ram}(L_{\mathrm{LY}}/\QQ).
\]
则
\[
\mathcal R_{10}\cap \mathcal R_{\mathrm{LY}}\subseteq \{3\}.
\]
因此，在任何有限 prime-register 或局部 bookkeeping 中，两侧可共享的有限局部通道至多只有一个 \(3\)-主方向；尤其 Lee--Yang 的 \(37\)-主局部缺陷对十次层 \(L_{10}\) 一侧严格不可见。
\end{theorem}
\begin{proof}
由定理 \ref{thm:xi-time-part9c-leyang-cubic-artin-weight1}，Lee--Yang 三次数域的分歧素数恰为
\[
\mathcal R_{\mathrm{LY}}=\{3,37\}.
\]
另一方面，定理 \ref{thm:xi-terminal-zm-delta-node-preimage-discriminant-and-norm} 给出
\[
\Disc(R_{10})=-2^{88}\cdot 3^{17}\cdot 11^{4}\cdot 31^{2}\cdot 727^{2}\cdot 5146068463^{2},
\]
并明确指出 \(37\) 不在判别式支撑中出现。由于分裂域的分歧只能出现在多项式判别式支撑所包含的素数上，故
\[
37\notin \mathcal R_{10}.
\]
于是
\[
\mathcal R_{10}\cap \mathcal R_{\mathrm{LY}}\subseteq \{3\}.
\]
最后一句只是这一包含关系在有限局部 bookkeeping 语言下的重述。证毕。
\end{proof}

\begin{theorem}[有限秩加性影子的统一障碍]\label{thm:conclusion-finite-rank-additive-shadow-obstruction}
\leanverified{paper\_conclusion\_finite\_rank\_additive\_shadow\_obstruction}
记
\[
\NN_{\times}:=(\NN_{\ge 1},\times).
\]
设 \(M\) 为任意有限生成交换可消去幺半群。则不存在幺半群单射
\[
\Psi:\NN_{\times}\hookrightarrow M.
\]
特别地，任何试图把乘法结构严格线性化为有限秩加法账本的方案，无论写成 \((\NN^k,+)\)、\((\ZZ\oplus L,+)\) 或其任意有限生成可消去变体，皆不可能。
\end{theorem}
\begin{proof}
这正是定理 \ref{thm:killo-prime-freedom-non-finitizability} 第 \(2\) 条的直接重述。若以 Grothendieck 完成表述，则定理 \ref{thm:killo-prime-freedom-non-finitizability} 第 \(1\) 条给出
\[
G(\NN_{\times})\cong \bigoplus_{p\in\mathbb P}\ZZ,
\]
其自由秩可数无穷；而任意有限生成交换可消去幺半群 \(M\) 的 Grothendieck 完成 \(G(M)\) 仅有有限 \(\ZZ\)-秩，故不可能容纳上述群的单射像。证毕。
\end{proof}

\begin{theorem}[有限秩精确账本完成类为空]\label{thm:conclusion-finite-rank-exact-ledger-completion-fiber-empty}
\leanverified{paper\_conclusion\_finite\_rank\_exact\_ledger\_completion\_fiber\_empty}
固定任意单射
\[
e:\NN_{\times}\hookrightarrow \ZZ.
\]
称一对 \((L,\Phi)\) 为 \(e\) 的有限秩精确账本完成，若 \(L\) 为有限生成交换群，且
\[
\Phi:\NN_{\times}\to \ZZ\oplus L
\]
是单射幺半群同态，并满足
\[
\mathrm{pr}_{\ZZ}\circ \Phi=e.
\]
则对每个这样的 \(e\)，其有限秩精确账本完成类都为空。等价地，遗忘函子
\[
U:\{\text{有限秩精确账本完成}\}\to \{\text{整数值单射编码}\},
\qquad
U(\Phi)=\mathrm{pr}_{\ZZ}\circ\Phi,
\]
在每个对象 \(e\) 上的纤维 \(U^{-1}(e)\) 都为空。
\end{theorem}
\begin{proof}
若 \(U^{-1}(e)\neq\varnothing\)，则存在有限生成交换群 \(L\) 与单射幺半群同态
\[
\Phi:\NN_{\times}\to \ZZ\oplus L.
\]
但这正被定理 \ref{thm:killo-godel-compression-not-finite-rank-homologizable} 所排除：只要 \(L\) 有限生成，就不存在这样的精确同态化账本。故每个纤维 \(U^{-1}(e)\) 必为空。

另一方面，定理 \ref{thm:conclusion-unbounded-prime-support-no-finite-register-embedding} 已说明，一旦转入可数素数估值向量
\[
n\longmapsto (v_p(n))_{p\in\mathbb P},
\]
便可获得忠实的幺半群表示。因而有限秩纤维为空并不意味着“无完成”，而是意味着一切精确完成都被迫跃迁到可数 prime-register 账本。证毕。
\end{proof}

\begin{corollary}[实现维与守恒维的极端分叉]\label{cor:conclusion-implementation-conservation-dimension-fork}
\leanverified{paper\_conclusion\_implementation\_conservation\_dimension\_fork}
沿用命题 \ref{prop:killo-implementation-vs-structural-half-circle-dimension} 的定义，
\[
\cdim^{\mathrm{impl}}(X):=\inf\left\{\frac{m}{2}:\ \exists\,\text{单射 }X\hookrightarrow \NN^m\right\},
\]
\[
\cdim^{\mathrm{hom}}_{+}(M):=\inf\left\{\frac{k}{2}:\ \exists\,\text{幺半群单射同态 }M\hookrightarrow (\NN^k,+)\right\}.
\]
则
\[
\cdim^{\mathrm{impl}}(\NN_{\times})=\frac12,
\qquad
\cdim^{\mathrm{hom}}_{+}(\NN_{\times})=+\infty.
\]
因而可实现性与可守恒性根本不是同一种维数问题：前者只要求集合嵌入，后者则要求乘法结构在有限秩加法账本上被严格保存。
\end{corollary}
\begin{proof}
命题 \ref{prop:killo-implementation-vs-structural-half-circle-dimension} 已给出这两条精确等式；其中第二式亦可由定理 \ref{thm:conclusion-finite-rank-additive-shadow-obstruction} 与定义直接推出。证毕。
\end{proof}

\paragraph{统一读法}
上述链条可压缩为：二维加权核中
\[
\Delta(z;u,v)=(1-vz^2)\mathcal Q_8(z;u,v)
\]
所分离出的参数盲刚性偶宇称支；bin-fold escort 温度族
\[
\{\pi_{m,q}\}_{q\ge 0}
\]
在极限中向单比特 Bernoulli 模型的二态闭包；六倍窗子序列
\[
m\equiv 4\pmod 6
\]
上被强制出现的半熵禁带以及碰撞/KL/峰值/最大纤维超额的同步压低；并且在最小三轴已审计窗口 \(m=58\) 上，这一半阶大陪集块又被推论 \ref{cor:conclusion-serrin-wulff-m58-odd-part-certificate} 锐化为一个带有非平凡 odd-part 的显式证书；十次层与 Lee--Yang 链之间
\[
\mathcal R_{10}\cap \mathcal R_{\mathrm{LY}}\subseteq \{3\}
\]
所刻画的唯一共同有限局部方向；以及
\[
\cdim^{\mathrm{impl}}(\NN_{\times})=\frac12,
\qquad
\cdim^{\mathrm{hom}}_{+}(\NN_{\times})=+\infty
\]
以及有限秩精确账本完成纤维的整体空性。它们共同表明：在本项目对象中，连续矩阵模态的可传播部分、escort 温度族的可辨识统计、六倍窗上强制湮灭的 Fourier 频带、可共享的有限局部分歧，以及可被有限加法账本忠实保存的算术结构，分别受彼此不同的刚性机制控制。统一性并不来自单一对称的塌缩，而来自这些机制在同一离散载体上的分层互补；一旦要求精确乘法守恒，有限秩账本的每条对象纤维都会整体坍缩为真空，只剩可数 prime-register 轴仍然可行。

\endinput
